\begin{textsample}{POS Dim 3 – 2016\_c6 – Score -5.00 – 2016\_c6/enade\_2016\_farm\_b3ee79edf9.txt}  \label{ex:f3_pos_255}
CAPÍTULO 3 – ANÁLISE TÉCNICA DA PROVA Este capítulo tem por objetivo apresentar o desempenho dos estudantes concluintes de Farmácia no Enade/2016.
Para isso, foram calculadas as estatísticas básicas da prova em seu todo ( seção 3.1.1 ), bem como as estatísticas relacionadas aos Componentes de Formação Geral ( seção 3.1.2 ) e de Conhecimento Específico da Área ( seção 3.1.3 ).
Dadas as suas características, foram analisadas, em separado, as questões objetivas ( seção 3.2 ) e as questões discursivas ( seção 3.3 ).
Para as questões objetivas, foram disponibilizados os índices de facilidade e de discriminação ponto-bisserial, também em separado, para os Componentes de Formação Geral ( seção 3.2.1 ) e de Conhecimento Específico ( seção 3.2.2 ).
De cada componente, uma das questões foi escolhida para exemplificar a análise gráfica, relacionando as alternativas escolhidas pelos estudantes ( inclusive o gabarito ) com o número de acertos no componente.
O Anexo I apresenta a íntegra da análise gráfica para todas as questões objetivas.
Para cada uma das questões discursivas, os conteúdos dos tipos mais comuns de respostas dos estudantes são apresentados e comparados com o padrão de respostas esperado ( ver Anexo VIII com o padrão de respostas ).
Tomando como base as duas questões discursivas do Componente de Formação Geral, a seção 3.3.1.6 apresenta comentários sobre a correção das respostas com respeito à Língua Portuguesa.
Nas Tabelas 3.1 e 3.3 são apresentados o tamanho da população inscrita, ausente, presente e percentual de faltosos, respectivamente, por Grande Região e por Categoria Administrativa e Organização Acadêmica da IES.
As demais tabelas apresentam as seguintes estatísticas das notas15 : média do desempenho na prova, erro padrão da média, desvio padrão, nota mínima, mediana e nota máxima para cada um de seus componentes.
Tais estatísticas contemplam o total de estudantes concluintes da Área de Farmácia inscritos e presentes à prova do Enade/2016, tendo em vista agregações, ou por Grandes Regiões e o país como um todo, ou por Categoria Administrativa e Organização Acadêmica da IES.
Em relação aos gráficos de distribuição de notas, o intervalo considerado foi de 10 unidades, aberto à esquerda e fechado à direita, com exceção do primeiro intervalo, [ 0 ; 10 ], fechado em ambos os extremos.
Para os gráficos de distribuição das notas das 15 Uma definição dessas estatísticas pode ser encontrada no Glossário. questão em branco16 e nota zero. 3.1 - CONDIÇÃO DE PRESENÇA E ESTATÍSTICAS BÁSICAS DA PROVA Esta seção apresenta a condição de presença, além das estatísticas selecionadas e histogramas da nota geral ( 3.1.1 ) e de cada componente : Formação Geral ( 3.1.2 ) e Conhecimento Específico ( 3.1.3 ).
São também apresentadas estatísticas selecionadas de subpopulações caracterizadas por Grande Região, Categoria Administrativa e Organização Acadêmica. 3.1.1 - Condição de Presença e Estatísticas Básicas Gerais A Tabela 3.1 apresenta a Condição de Presença na prova, por Grande Região, dos estudantes concluintes de Farmácia.
A população total de Inscritos foi de 13.991.
Destes, 13.197 estiveram Presentes, sendo 5,7 \% o índice de não comparecimento.
A Região de maior ausência foi a Sudeste ( 6,8 \%, ) e a de menor ausência foi a Sul ( 3,7 \% ).
Já a Tabela 3.2 apresenta as Estatísticas Básicas da prova, por Grande Região, dos estudantes concluintes de Farmácia.
A Média das notas da prova como um todo ( nas seções seguintes serão analisados os Componentes de Formação Geral e de Conhecimento Específico ) foi 50,5, sendo que os alunos da região Norte obtiveram a Média mais baixa ( 43,6 ), e os da região Sul obtiveram a Média mais alta ( 53,6 ).
As demais médias foram : 49,8 na região Nordeste, 51,9 na região Sudeste e 48,7 na região Centro-Oeste.
O Desvio padrão para o Brasil como um todo foi 13,5, sendo o maior Desvio padrão encontrado na região Centro-Oeste ( 13,6 ) e o menor, na região Sudeste ( 13,1 ), indicando que a variação da dispersão foi pequena entre as regiões.
A região que obteve a maior nota Máxima foi a Sudeste ( 92,0 ), ao passo que a região que atingiu a menor nota Máxima foi a Norte ( 80,5 ).
A Mediana do Brasil como um todo foi 50,7, sendo a maior Mediana obtida na região Sul ( 54,3 ), e a menor obtida na Norte ( 43,1 ).
A maior nota Mínima ocorreu na região Nordeste ( 8,7 ).
Na região Norte a nota Mínima foi 5,6, na região Sul foi 5,3 e nas demais regiões foi zero.
Considerando -se as notas segundo Grande Região, observa -se que existe diferença estatisticamente significativa ao nível de 95 \% entre a menor Média, obtida na região Norte ( 43,6 ), e as médias de todas as demais regiões.17 O comportamento das notas dos estudantes de todo o Brasil pode ser observado no Gráfico 3.1 que apresenta um histograma com a distribuição das mesmas.
Essa é uma distribuição unimodal com a moda no intervalo ( 50 ; 60 ], seguido bem de perto pelo intervalo ( 40 ; 50 ]. 17 Todas as comparações deste capítulo utilizam os intervalos de 95 \%.
Os erros-padrão da média que possibilitam os testes estão disponíveis nas tabelas.
A Tabela 3.3 apresenta informações referentes à Condição de Presença desagregadas por Categoria Administrativa e por Organização Acadêmica.
Da população total de Inscritos, 10.311 são de IES Privadas e 3.680 de IES Públicas.
Em relação à Organização Acadêmica, a maior participação foi obtida por estudantes de Universidades ( 6.939 ), seguida pelos de Faculdades ( 4.586 ) e Centros Universitários ( 2.466 ).
A Categoria Administrativa de maior ausência foi a Privada ( 6,0 \% ), e entre as Organizações Acadêmicas foi a dos Centros Universitários ( 7,6 \% ), que ficou acima da taxa de ausência nacional de 5,7 \%.
Já a Tabela 3.4 apresenta informações referentes à Média da nota final desagregadas por Categoria Administrativa e por Organização Acadêmica.
A Média das notas da prova como um todo foi 50,5.
Em relação à Categoria Administrativa, os estudantes das IES Públicas obtiveram Média mais alta ( 57,3 ), e os das IES Privadas obtiveram Média mais baixa ( 48,0 ) que a Média nacional.
Observa -se que existe diferença estatisticamente significativa entre as médias das notas das IES Públicas e Privadas.
A diferença entre as médias das regiões Sul e Norte ( 10,0 ), a maior e a menor Média, é superior à diferença entre IES Públicas e Privadas ( 9,3 ), caracterizando uma diversidade administrativa um pouco menor do que regional.
O Desvio padrão para as IES Públicas ( 12,5 ) foi inferior ao do Brasil como um todo ( 13,5 ), indicando uma dispersão um pouco menor das notas nesta Categoria Administrativa.
No tocante à Organização Acadêmica, as Universidades obtiveram Média mais alta que a nacional ( 54,5 e 50,5, respectivamente ).
A média das Faculdades foi 45,8.
Constata -se que existe diferença estatisticamente significativa ao nível de 95 \% entre a Média das Universidades e dos outros dois tipos de Organização Acadêmica. 3.1.2 - Estatísticas Básicas no Componente de Formação Geral A Tabela 3.5 apresenta as Estatísticas Básicas em relação ao componente da prova que avalia a Formação Geral dos estudantes concluintes.
Os alunos de todo o Brasil obtiveram desempenho médio de 46,6.
Quanto à variabilidade, o Desvio padrão das notas dos estudantes do Brasil, como um todo, foi 16,5.
A maior Média foi obtida na região Sudeste ( 48,4 ), e a menor, na região Norte ( 41,2 ).
As demais médias foram : 45,3 na região Nordeste, 48,2 na região Sul e 44,6 na região Centro-Oeste.
Já o maior Desvio padrão foi obtido na região Norte ( 16,8 ), e o menor, na região Sudeste ( 16,1 ).
Os demais desvios padrões foram : 16,4 na região Nordeste, 16,5 na região Sul e 16,3 da região Centro-Oeste.
A maior nota no Componente de Formação Geral da prova do Enade/2016 foi 95,8, obtida por, pelo menos, um aluno na região Sul.
A menor nota Máxima foi obtida na região Centro-Oeste ( 91,2 ).
A Mediana do Brasil, como um todo, foi 47,1, sendo a menor Mediana encontrada na região Norte ( 41,0 ), e a maior, encontrada na região Sudeste ( 48,8 ).
A nota Mínima nesta parte foi zero em todas as regiões, sem exceção.
Considerando -se as notas segundo Grande Região, observa -se que existe diferença estatisticamente significativa entre a maior Média das notas do Componente de Formação Geral, obtida na região Sudeste ( 48,4 ), e a Média das regiões Norte, Nordeste e Centro-Oeste.
No entanto, não há diferença estatisticamente significativa entre as médias das regiões Sudeste ( 48,4 ) e Sul ( 48,2 ).
O Gráfico 3.2 propicia a avaliação do desempenho dos estudantes do Componente de Formação Geral, a partir do histograma da distribuição das notas correspondentes.
A distribuição é unimodal, com moda em ( 40 ; 50 ], enquanto na prova, como um todo ( Gráfico 3.1 ), a moda foi alcançada no intervalo ( 50 ; 60 ].
Nota -se, ainda que, no Gráfico 3.2, as notas apresentam maior dispersão do que no Gráfico 3.1 ( distribuição das notas da prova ), confirmado pela comparação dos desvios padrões : 13,5 para a nota da prova como um todo e 16,5 para o Componente de Formação Geral.
Na Tabela 3.6 são apresentadas as informações referentes ao desempenho dos concluintes do Componente de Formação Geral, em diferentes agregações : Categoria Administrativa e Organização Acadêmica.
Observa -se que existe diferença estatisticamente significativa entre as médias dos tipos de Categoria Administrativa.
A maior Média foi obtida por estudantes de IES Públicas ( 53,2 ), com uma diferença estatisticamente significativa da obtida por estudantes de IES Privadas ( 44,1 ).
Nota -se que também há diferença estatisticamente significativa entre as médias dos três tipos de Organização Acadêmica, valor maior para Universidades e menor para Faculdades. 3.1.3 - Estatísticas Básicas do Componente de Conhecimento Específico A Tabela 3.7 apresenta as Estatísticas Básicas referentes ao Componente de Conhecimento Específico da área de Farmácia.
A Média do desempenho dos alunos do Brasil, como um todo, foi 51,8.
A maior Média foi obtida na região Sul ( 55,5 ), e a menor, na região Norte ( 44,4 ).
As demais médias foram : 51,3 na região Nordeste, 53,0 na região Sudeste e 50,0 na região Centro-Oeste.
Quanto à variabilidade das notas, o Desvio padrão do Brasil, como um todo, foi 14,8, sendo o maior Desvio padrão observado na região Centro-Oeste ( 14,9 ), e o menor, na região Norte ( 14,3 ).
Os demais desvios foram : 14,7 na região Nordeste, 14,4 na região Sudeste e 14,5 na região Sul.
A Mediana das notas dos estudantes de todo o Brasil foi 52,0.
A maior Mediana ocorreu na região Sul ( 56,2 ), e a menor, na região Norte ( 44,3 ).
As demais medianas foram : 51,3 na região Nordeste, 53,3 na região Sudeste e 50,2 na região Centro-Oeste.
A nota Máxima do Brasil, como um todo, foi 99,0, sendo obtida por, pelo menos, um aluno na região Sudeste.
As demais notas máximas foram : 87,1 na região Norte, 94,0 na região Nordeste, 95,5 na região Sul e 97,0 na região Centro-Oeste.
A nota Mínima foi diferente de zero nas regiões Nordeste ( 3,5 ) e Sul ( 4,5 ).
Observa -se que existe diferença estatisticamente significativa da menor Média das notas do Componente de Conhecimento Específico da região Norte ( 44,4 ) em relação às demais regiões.
Assim como os Gráficos 3.1 e 3.2, o Gráfico 3.3, apresentado a seguir, permite uma avaliação do desempenho de concluintes de Farmácia em relação ao Componente de Conhecimento Específico com um histograma da distribuição das notas correspondentes.
Esta também é uma distribuição unimodal, e o grupo modal é o ( 50 ; 60 ], o mesmo grupo modal da prova como um todo e acima do grupo modal para a Formação Geral.
A Tabela 3.8 apresenta uma comparação dos resultados em relação à Categoria Administrativa e à Organização Acadêmica, agora levando em conta o desempenho de alunos do Componente de Conhecimento Específico da prova.
No que se refere à Organização Acadêmica, a maior Média foi das Universidades ( 55,9 ), vindo a seguir a dos Centros Universitários ( 49,3 ) e, depois, a das Faculdades ( 46,9 ).
O maior Desvio padrão, e acima do valor para o Brasil, como um todo, foi o das Universidades ( 14,4 ).
As Universidades também obtiveram as maiores notas Máxima ( 99,0 ) e Mediana ( 56,2 ).
Os Centros Universitários obtiveram nota Máxima igual a 88,1 e as Faculdades obtiveram nota Máxima igual a 93,7.
As medianas foram : 49,5 nos Centros Universitários e 46,8 nas Faculdades.
A nota Mínima foi igual a zero para Universidades e Faculdades e 1,5 nos Centros Universitários.
Observa -se que existem diferenças estatisticamente significativas ao nível de 95 \% no Componente de Conhecimento Específico entre as notas das Organizações Acadêmicas.
Quanto à Categoria Administrativa, observa -se um comportamento semelhante àquele da parte de Formação Geral e à prova como um todo, ou seja, existe diferença estatisticamente significativa entre as médias das IES Públicas ( 58,6 ) e IES Privadas ( 49,3 ).
Neste caso também, a maior Média foi obtida por alunos de IES Públicas de ensino. 3.2 - ANÁLISE DAS QUESTÕES OBJETIVAS Esta seção apresenta estatísticas selecionadas e histogramas das Questões Objetivas de Formação Geral ( 3.2.1 ).
São também apresentadas e comparadas as médias das subpopulações caracterizadas por Grande Região. 3.2.1 - Componente de Formação Geral A Tabela 3.9 apresenta as Estatísticas Básicas relativas às oito questões objetivas do componente da prova que abrange a Formação Geral dos estudantes.
A Média do Brasil foi 41,0.
A menor Média foi encontrada na região Norte ( 35,4 ), e a maior, na região Sul ( 43,0 ).
As demais médias foram : 39,9 na região Nordeste, 42,9 na região Sudeste e 38,5 na região Centro-Oeste.
O Desvio padrão do Brasil foi 20,4, sendo o maior Desvio padrão encontrado na região Sudeste ( 20,5 ), e o menor, na região Norte ( 19,5 ).
Os demais desvios foram : 20,1 na região Nordeste, 20,2 na região Sul e 20,0 na região Centro-Oeste.
As medianas do Brasil como um todo e a das cinco regiões foram 37,5.
As notas máximas ( 100,0 ) e as notas mínimas ( 0,0 ) foram iguais para todas as regiões.
A Tabela 3.10 apresenta o Índice de Facilidade e o Índice de Discriminação ( ponto-bisserial ) para cada uma das questões objetivas do Componente de Formação Geral.
Quanto ao Índice de Facilidade, foram usadas as seguintes cores para diferenciar o nível de dificuldade da questão : Já quanto ao Índice de Discriminação, foram usadas as seguintes cores para qualificar a questão : As questões objetivas do Componente de Formação Geral, segundo o Índice de Facilidade, foram assim avaliadas : das oito questões, nenhuma teve o Índice de Facilidade classificado como Muito fácil e uma questão foi tida como Fácil, por ter índice de acertos 0,79.
Três questões foram consideradas com índice de dificuldade Médio, situando -se no intervalo entre 0,41 e 0,60 do Índice de Facilidade, ou seja, houve entre 41,0 \% e 60,0 \% de acertos, também foram três as questões classificadas na categoria Difícil, situando -se no intervalo entre 0,16 e 0,40.
Por fim, uma questão apresentou menos de 15 \% de acertos ( 12 \% ), razão pela qual seria classificada como Muito difícil.
Como já comentado, para análise das questões objetivas, relativas à Formação Geral, segundo o poder de discriminação, utilizou -se o Índice de Discriminação ponto-bisserial.
Nesta análise, as questões foram assim avaliadas : sete das oito questões apresentaram índice acima ou igual a 0,40 e, assim, foram classificadas com índice Muito bom para esse grupo de estudantes.
Uma questão teve Índice de Discriminação Bom, com valor entre 0,30 e 0,39.
Nenhuma questão teve nível Médio ou Fraco de discriminação para esse grupo de estudantes.
O Índice de Facilidade variou de 0,12 a 0,79, e o de Discriminação, de 0,30 a 0,55.
As sete questões com Índice de Discriminação Muito bom figuraram entre níveis de dificuldade Fácil, Médio e Difícil : uma classificada na categoria Fácil ( a questões 6 ) do Índice de Facilidade, três, na categoria Médio ( questões 1, 2 e 3 ) e três na categoria Difícil ( questões 4, 7 e 8 ).
Em particular, a questão 3 foi a que apresentou o maior poder discriminatório, com índice 0,55, e foi considerada de nível Médio em termos de facilidade, com uma proporção de 0,60 acertos.
O máximo de acertos foi alcançado pela questão 6 com um Índice de Facilidade de 0,79.
A questão de número 4 apresentou um Índice de Facilidade de 0,24, ou seja, um pouco menos do que um quarto dos estudantes conseguiu resolvê -la, dentro do universo de participantes.
Seu Índice de Discriminação foi Muito Bom ( 0,46 ).
Já a questão 5 obteve Índice de Discriminação Bom, 0,30, e seu Índice de Facilidade foi Muito Difícil ( 0,12 ).
O Gráfico 3.4, para exemplificar, analisa o comportamento da questão de número 2 de Formação Geral.
Trata -se de uma questão com índice de facilidade Médio e a que obteve o segundo maior Índice de Discriminação dessa parte da prova.
Neste gráfico, cada uma das cinco curvas representa o percentual de respostas em determinada alternativa da questão, em função do número de acertos dos estudantes nessa parte da prova ( Formação Geral/Múltipla Escolha ), antes de possíveis eliminações pelo critério do ponto-bisserial.
A curva em azul corresponde à alternativa A, a alternativa correta para esta questão.
Observa -se que entre os estudantes com até dois acertos, nessa parte do exame, a situação mais frequente foi a escolha da alternativa E ( em vermelho ), incorreta.
Por exemplo, entre os estudantes que acertaram três questões, 25,7 \% escolheram a alternativa E, 29,9 \% escolheram a alternativa B ( em verde ), 2,1 \% escolheram a alternativa C ( preto ) e 1,3 \% a D ( roxo ), enquanto quase metade ( 40,8 \% ) escolheu o gabarito, letra A.
Entre os que acertaram três respostas das questões de múltipla escolha de Formação Geral, 0,1 \% deixou a questão em branco e 0,1 \% marcou mais de uma alternativa, invalidando a questão.
À medida que o número de acertos aumenta, indicando desempenho melhor nessa parte da prova, aumenta concomitantemente a proporção de estudantes que selecionaram a alternativa correta A, atingindo 100 \% para os estudantes com 8 acertos.
Essa análise permite verificar como a questão discriminou os grupos de desempenho, justificando o alto índice obtido na questão.
Cumpre notar que não é possível inferir deste gráfico nem o índice de facilidade ( que seria uma média da proporção ponderada pela quantidade de alunos com cada uma das notas ), nem o índice de discriminação ponto-bisserial, por razão equivalente.
No caso extremo no qual a grande concentração dos acertos dos alunos fosse abaixo de 4, o índice de facilidade seria obrigatoriamente abaixo de 40 \% ( neste exemplo ).
Caso a concentração fosse em 6 acertos ou mais, o índice seria obrigatoriamente acima de 60 \%.
Os gráficos relativos às demais questões de Formação Geral constam do Anexo I. 3.2.2 - Componente de Conhecimento Específico A Tabela 3.11 apresenta as Estatísticas Básicas em relação às questões objetivas do Componente de Conhecimento Específico da prova de Farmácia por Grande Região.
A Média do Brasil deste componente foi de 51,2.
A menor Média foi observada na região Norte ( 44,2 ), e a maior, na região Sul ( 55,0 ).
O Desvio padrão de todo o Brasil foi 15,8, sendo o menor Desvio padrão encontrado na região Norte ( 15,1 ), e o maior, na região Centro-Oeste ( 15,9 ).
A Mediana de todo o Brasil foi 52,6, a mesma encontrada em duas das regiões : Nordeste e Sudeste.
Nas regiões Norte e Centro-Oeste a Mediana foi menor ( 42,1 e 47,4, respectivamente ) e na região Sul, maior ( 57,9 ).
A nota Máxima da prova foi 100, obtida nas questões objetivas do Componente de Conhecimento Específico, por, pelo A Tabela 3.12 apresenta os Índices de Facilidade e Discriminação ( ponto-bisserial ) das questões objetivas do Componente de Conhecimento Específico para os estudantes de Farmácia.
Para facilitar a diferenciação das questões, usaremos as mesmas cores da Tabela 3.10 para as diferentes classificações dos Índices de Facilidade e de Discriminação.
A partir do Índice de Facilidade obtido, pode -se concluir que 44 \% das questões objetivas da prova foi considerado, pelo menos, Difícil : das 27 questões válidas, 12 foram classificadas como Difícil ou como Muito difícil.
Nove questões foram classificadas como Fácil, e outras seis consideradas como Médio.
Já quanto aos índices de discriminação das questões objetivas do Componente de Conhecimento Específico da prova, tem -se como resultado a seguinte classificação : 11 das 27 questões foram consideradas como boas, enquanto uma delas obteve Índice de Discriminação Muito bom.
Assim, para um pouco menos do que a metade das questões – 12 em 27 – os Índices de Discriminação foram Bom ou Muito bom.
Dentre as demais, sete delas foram classificadas como Médio e outras oito como Fraco, sendo 15, por conseguinte, a quantidade de questões nos dois patamares mais baixos de discriminação.
Constata -se, assim, que a prova – no que se refere ao Componente de Conhecimento Específico – possuía capacidade razoável de discriminar entre aqueles que dominam ou não o conteúdo.
Dentre as questões que alcançaram os maiores índices de discriminação, uma delas, a de número 27, foi classificada com Índice de Discriminação Muito bom, situando -se no intervalo de 0,40 a 0,50 do índice, e foi classificada como Fácil, quanto ao Índice de Facilidade.
A questão de número 30 foi a mais difícil dentre as 27 questões menos, um aluno nas regiões Sudeste, Sul e Centro-Oeste.
Na região Norte a nota Máxima da prova foi 89,5 e na região Nordeste foi 94,7.
A única região onde a nota Mínima foi diferente de zero foi a Sul ( 5,3 ).
A Tabela 3.12 apresenta os Índices de Facilidade e Discriminação ( ponto-bisserial ) das questões objetivas do Componente de Conhecimento Específico para os estudantes de Farmácia.
Para facilitar a diferenciação das questões, usaremos as mesmas cores da Tabela 3.10 para as diferentes classificações dos Índices de Facilidade e de Discriminação.
A partir do Índice de Facilidade obtido, pode -se concluir que 44 \% das questões objetivas da prova foi considerado, pelo menos, Difícil : das 27 questões válidas, 12 foram classificadas como Difícil ou como Muito difícil.
Nove questões foram classificadas como Fácil, e outras seis consideradas como Médio.
Já quanto aos índices de discriminação das questões objetivas do Componente de Conhecimento Específico da prova, tem -se como resultado a seguinte classificação : 11 das 27 questões foram consideradas como boas, enquanto uma delas obteve Índice de Discriminação Muito bom.
Assim, para um pouco menos do que a metade das questões – 12 em 27 – os Índices de Discriminação foram Bom ou Muito bom.
Dentre as demais, sete delas foram classificadas como Médio e outras oito como Fraco, sendo 15, por conseguinte, a quantidade de questões nos dois patamares mais baixos de discriminação.
Constata -se, assim, que a prova – no que se refere ao Componente de Conhecimento Específico – possuía capacidade razoável de discriminar entre aqueles que dominam ou não o conteúdo.
Dentre as questões que alcançaram os maiores índices de discriminação, uma delas, a de número 27, foi classificada com Índice de Discriminação Muito bom, situando -se no intervalo de 0,40 a 0,50 do índice, e foi classificada como Fácil, quanto ao Índice de Facilidade.
A questão de número 30 foi a mais difícil dentre as 27 questões específicas, com baixo Índice de Facilidade, apenas 7,0 \% de acertos.
Essa questão apresentou poder discriminatório igualmente baixo, 0,04, o que comprova ter sido esta questão a mais difícil para os estudantes.
Destacam -se, também, as questões 12 e 18, com Índice de Facilidade 0,08 e 0,15, respectivamente, o que, em termos percentuais, corresponde a 8 \% e 15 \% de estudantes terem respondido acertadamente.
Já o Índice de Discriminação dessas questões foi 0,15 ( questão 12 ) e 0,13 ( questão 18 ).
Tais questões foram, portanto, pelo critério ponto-bisserial, consideradas inadequadas.
Por isso, as questões 12,18 e 30, além das questões 1, 14, 17, 24 e 31 foram eliminadas do cômputo da nota final.
A título de exemplo das análises do comportamento das questões objetivas, o Gráfico 3.5 analisa a questão 27 do Componente de Conhecimento Específico.
Esta questão foi considerada pelas respostas dos estudantes avaliados na prova como uma questão com nível de facilidade Fácil, 0,65, ou seja, 65,0 \% dos estudantes assinalaram acertadamente a opção D, correspondente ao gabarito.
Seu Índice de Discriminação foi igual a 0,47, classificado como Muito bom, também sendo esta questão a que apresentou o maior valor.
Neste gráfico, cada uma das cinco curvas representa o percentual de respostas em determinada alternativa da questão 27, em função do número de acertos dos estudantes nessa parte da prova, antes de possíveis eliminações de questões pelo critério do ponto-bisserial.
A alternativa correta D, representada no gráfico pela curva em roxo, foi escolhida em maiores proporções pelos alunos com desempenho melhor nessa parte da prova.
Já as alternativas incorretas, também denominadas distratores, foram selecionadas, principalmente, por aqueles com notas mais baixas.
Neste caso também a soma não é sempre 100 \% por causa das questões não respondidas ou com mais de uma opção marcada.
Aqueles com nota zero, na sua quase totalidade deixaram esta questão em branco ou marcaram mais de uma alternativa, comportamento considerado inválido.
A proporção de alunos que selecionaram a resposta correta D aumenta gradativamente, chegando a atingir 100 \% para os que acertaram 20 ou mais questões, enquanto a proporção dos que escolheram alternativas incorretas decai, a partir de quatro acertos, como função do número de acertos nessa parte da prova.
Os gráficos relativos às demais questões do Conhecimento Específico constam do Anexo I. 3.3 - ANÁLISE DAS QUESTÕES DISCURSIVAS Esta seção apresenta estatísticas selecionadas e histogramas das Questões Discursivas de Formação Geral ( 3.3.1 ) e Conhecimento Específico ( 3.3.2 ).
São também apresentadas e comparadas as médias de subpopulações caracterizadas por Grande Região. 3.3.1 - Componente de Formação Geral As análises dos resultados de desempenho dos estudantes de Farmácia nas duas questões discursivas relativas à Formação Geral encontram -se na Tabela 3.13 e no Gráfico 3.6.
Na Tabela 3.13, observa -se que a nota Média nesse conjunto de questões foi acima da obtida nas objetivas.
Os estudantes, de todo o Brasil, obtiveram, em Formação Geral, Média 41,0 nas questões objetivas e 54,8 nas questões discursivas.
Pode -se notar também que o Desvio padrão nesse conjunto de questões foi um pouco maior do que o obtido nas objetivas, 20,4 nas questões objetivas e 21,6 nas questões discursivas.
A maior Média foi obtida na região Sudeste ( 56,6 ), e a menor, na região Norte ( 50,0 ).
A Mediana de todo o Brasil, neste componente, foi 60,0, a mesma obtida na região Centro-Oeste.
Nas regiões Norte e Nordeste a Mediana foi menor, 56,0 e 59,0, respectivamente.
Nas regiões Sudeste e Sul a Mediana foi 61,5, maior do que a do Brasil como um todo.
A nota Máxima do Brasil como um todo ( 97,0 ) foi obtida na região Nordeste.
Nas demais regiões a nota Máxima foi : 92,5 na região Norte, 96,0 na região Sudeste, 95,5 na região Sul e 95,0 na região Centro-Oeste.
A nota Mínima ( 0,0 ) foi a mesma em todas as regiões do Brasil.
O Gráfico 3.6 representa a distribuição das notas nas questões discursivas do Componente de Formação Geral.
A moda desta distribuição ocorre no intervalo ( 60 ; 70 ].
Destaca -se também o intervalo [ 0 ; 10 ] com distribuição próxima a 8 \% do total de notas, sendo que no intervalo [ 0 ; 10 ] incluem -se além da nota zero, a frequência de alunos que deixaram este tipo de questão em branco.
Na sequência, os resultados verificados para cada uma das questões discursivas de Formação Geral serão apresentados, estabelecendo -se relações com os conteúdos abordados em cada uma delas.
Os comentários da Banca de docentes corretores, a respeito do observado na correção das respostas dos estudantes, suas impressões e conclusões serão apresentados junto à análise de cada questão.
Cumpre esclarecer que, tendo em vista que as questões discursivas de Formação Geral são padronizadas, ou seja, constam de todas as provas, os comentários da Banca são os mesmos para todas as carreiras acadêmicas, sendo direcionados a todos os estudantes que participaram do Enade/2016.
A seguir, serão analisados os desempenhos da Área de Farmácia nas duas questões discursivas de Formação Geral do Enade/2016, comparando -se os resultados obtidos com comentários para cada questão. 3.3.1.1 - Análise de Conteúdo da Questão Discursiva 1 do Componente de Formação Geral Os dados de Farmácia, obtidos a partir das respostas à questão 1, encontram -se na Tabela 3.14 e no Gráfico 3.7.
Nessa questão – de melhor desempenho dentre as duas de Formação Geral – os alunos, de todo o Brasil, tiveram Média 61,1.
A maior Média para a questão 1 foi obtida na região Sudeste ( 63,3 ), e a menor, na região Norte ( 54,3 ).
Quanto à variabilidade das notas, o Desvio padrão de todo o Brasil foi 30,6.
O menor Desvio padrão foi obtido na região Sudeste ( 29,4 ) e o maior Desvio padrão foi obtido na região Norte ( 32,9 ).
As medianas de quatro das cinco regiões e a do Brasil como um todo foram iguais ( 70,0 ), a exceção foi a região Norte com mediana 60,0.
As notas máximas e mínimas da questão discursiva 1 foram, sem exceção, as mesmas para todas as regiões do Brasil, respectivamente, 100,0 e 0,0.
O Gráfico 3.7 mostra a distribuição das notas na questão discursiva 1 do Componente de Formação Geral.
Observa -se que a maior frequência corresponde aos alunos que obtiveram nota dentro do intervalo ( 70 ; 80 ], correspondendo à moda da distribuição.
Destacam -se, também, os alunos que deixaram a questão em branco, com aproximadamente 12 \% do total. 3.3.1.2 - Comentários sobre a correção de Conteúdo das respostas à Questão Discursiva 1 A questão discursiva 1 alinha -se aos conteúdos programáticos previstos nas diretrizes curriculares, quanto às habilidades e às competências para a atualização permanente e aos conhecimentos sobre a realidade brasileira, mundial, tangenciando outras áreas do conhecimento.
O enunciado é claro e indica um tema de extrema relevância, qual seja, as migrações humanas e todo o complexo de problemas que esse movimento causa, quer na origem, quer no seu ponto de destino.
A questão é extremamente atual, maximizada pelos acompanhamentos midiáticos que espetacularizam cenas dramáticas do processo de fuga de massas humanas das guerras, como a morte de imigrantes.
Além do aspecto internacional, deve se ter em conta que o Brasil sofreu vários influxos migratórios, desde o forçado, pela transferência de escravos da África, até o incentivado, dos italianos, alemães e outros europeus para trabalhar na agricultura.
Isso sem considerar a própria colonização portuguesa originária e as migrações posteriores dos mesmos portugueses.
Posteriormente, o Brasil tornou -se um país de emigrantes que foram aportar a outros locais na busca de melhor situação econômica.
Muitos migraram para a América do Norte e para a Europa, principalmente.
Mais recentemente, por força do crescimento econômico pátrio, novas levas de imigrantes desembarcaram no Brasil, muitos deles provenientes do Haiti, embora tenhamos recebidos vizinhos latinos, como bolivianos, paraguaios, chilenos, argentinos, venezuelanos, peruanos, colombianos que também vieram em busca de uma vida melhor.
Como se verifica, a temática indicada no enunciado pertence à realidade conhecida amplamente pelos estudantes, pois, também divulgada por todos os meios de comunicação.
Trata -se de tema geral que, normalmente, não se inclui nas grades curriculares superiores, mas é de conhecimento geral e relacionado aos direitos humanos, este sim, ponto de integração às diretrizes curriculares.
A linguagem utilizada foi clara, indicando a existência dos fatos vinculados aos movimentos migratórios.
Por outro lado, o texto demanda uma reflexão sobre a realidade dos indivíduos que são compelidos ao afastamento dos seus lares por motivos externos, como guerras, conflitos civis, religiosos, perseguições das mais variadas, questões econômicas.
Como se trata de texto contendo uma ampla gama de possibilidades, é natural que ocorram interpretações diversas das esperadas.
Do mesmo modo, a contemporaneidade do tema encaminha o enunciado para ser qualificado como questão de média dificuldade, em comparação com as demais questões apresentadas.
O padrão de resposta indicou possíveis interpretações do texto.
No que concerne ao grau de profundidade, mostrou a necessidade de conhecimento de vários itens pertinentes à migração e que são refletidos nas várias possibilidades de resposta previstas pelo padrão.
Pela análise das provas realizadas, constatou -se que muitos foram os que identificaram os itens indicados no padrão de resposta.
Essa circunstância poderia ser justificada, tendo em vista que o tema vem sendo bastante difundido na mídia social, gerando elementos naturais para a resposta, bem como da história do país, que, intimamente, vincula -se ao processo migratório, quer ativo, remetendo pessoas para outros países, quer passivo, na recepção de pessoas.
Essa característica também permite responder aos termos postos quanto aos Estados que recebem os migrantes.
Assim, a indicação dos aspectos positivos e negativos também é bastante difundida.
As respostas medianas e boas apreenderam o sentido do enunciado, e indicaram os motivos expressos no padrão de respostas, como os ensejadores das imigrações, notadamente, para a Europa, que seria o foco do texto.
Por outro lado, também identificaram os problemas dos imigrantes, no mesmo padrão indicado para as respostas desejadas.
Alguns concluintes procuraram exemplos de comparação na realidade brasileira, como os clássicos imigrantes italianos, alemães e de outras nacionalidades, bem como os atuais haitianos e demais povos da América Latina.
Uma minoria não apreendeu o sentido do enunciado e, somente, apresentou elementos da realidade nacional.
Esses foram avaliados no conjunto de provas fracas.
As respostas demonstraram que o concluinte está atualizado, em regra, com o tema da imigração, indicando os fundamentos dessa movimentação populacional que, nos últimos tempos, tem se agravado pela incidência de conflitos armados em vários países, ensejadores de reflexos na segurança dos habitantes desses locais, bem como o surgimento de focos de doenças e ausência dos serviços públicos básicos, além da fome, decorrente da impossibilidade de produzir alimentos ou mesmo comerciá -los.
Do ponto de vista dos Estados que recebem o fluxo de imigrantes, também restou clara a percepção das dificuldades vinculadas à ausência de estrutura quanto aos serviços públicos ( saúde, educação, segurança, moradia, alimentação ) para receber um número expressivo de imigrantes.
Em época de recessão econômica, a presença de estrangeiros acarreta a competição pela busca de empregos, gera movimentos xenofóbicos, bem como choques culturais e políticos.
Deve -se ter em mira que esses movimentos têm um viés de insegurança, pelos países receptores, diante da crise causada por movimentos considerados terroristas, embora a imensa maioria não tenha qualquer vínculo com esses movimentos.
Os recentes atentados na Europa indicam que os seus autores têm origem em comunidades de imigrantes, embora muitos já tenham nascido nos países da comunidade europeia.
Essas características negativas, no entanto, paradoxalmente, convivem com aspectos positivos, o que encaminha a questão para a necessidade dos jovens imigrantes, que ocuparão postos de trabalho não desejados pelos habitantes locais, bem como, compensarão os baixos índices de crescimento populacional.
Existe outro reflexo positivo que está relacionado ao sistema previdenciário, que terá novos contribuintes nos trabalhadores estrangeiros que se integram ao sistema.
Esses pontos foram abordados e indicam que existe uma consciência dos concluintes dos cursos universitários quanto aos principais itens relacionados ao tema que foi posto para análise. 3.3.1.3 - Análise de Conteúdo da Questão Discursiva 2 do Componente de Formação Geral A Tabela 3.15 mostra que o desempenho médio dos estudantes na questão discursiva 2 ( média 46,2 ) foi inferior ao obtido na questão discursiva 1 ( média 61,1 ).
A região Sudeste foi aquela cuja Média, também nessa questão, foi maior ( 47,5 ), e a de menor Média foi a região Norte ( 43,6 ).
Quanto à variabilidade das notas, o Desvio padrão de todo o Brasil foi 24,2, inferior ao obtido na questão discursiva 1 ( 30,6 ).
O maior desvio nessa questão foi obtido na região Norte ( 25,0 ), enquanto o menor foi obtido na região Sudeste ( 23,8 ).
A Mediana de todo o Brasil foi 50,0, a mesma de todas as regiões.
As notas máximas ( 100,0 ) e as notas mínimas ( 0,0 ) também foram as mesmas em todas as regiões, sem exceção.
O Gráfico 3.8 mostra a distribuição das notas na questão discursiva 2 do Componente de Formação Geral.
Observa -se que a maior frequência corresponde aos alunos que obtiveram nota dentro do intervalo ( 50 ; 60 ].
Destacam -se, também, os alunos que deixaram a questão em branco, caracterizando -se como máximo local.
Nota -se, ainda, que as notas ficaram mais concentradas em comparação à questão discursiva de número 1.
Isto pode ser constatado também pela comparação do desvio padrão das notas da questão discursiva 2 ( 24,2 ) e da questão discursiva 1 ( 30,6 ). 3.3.1.4 - Comentários sobre a correção de Conteúdo das respostas à Questão Discursiva 2 A questão discursiva 2 tem seu conteúdo integrado às Diretrizes Curriculares, uma vez que a proteção aos direitos humanos e às relações de gênero constituem matrizes fundamentais para conhecimento dos profissionais de nível superior.
O enunciado é claro e indica um tema de extrema relevância, qual seja, a análise da questão da mulher submetida a injustas agressões e os seus reflexos pessoais, sociais e familiares.
Também enfrenta a busca pela maior proteção no concernente às agressões, com a instituição de um modelo legal mais repressivo e que foi gerado por uma situação especifica de ataque à mulher, com consequências dramáticas para a vítima, Maria da Penha, cuja lei levou seu nome.
Deve ser recordado que os conflitos vinculados às relações familiares sempre foram circunscritos ao seu âmbito e, quando dele extravasavam, por força de agressões contundentes à integridade física de um dos partícipes da relação, em geral a mulher, raramente redundavam em punições relevantes, sendo comum o pagamento em cestas básicas de alimentos, como sanção.
Após a edição da lei Maria da Penha ( Lei no. 11.340, sancionada em 7 de agosto de 2006 ), além das sanções sofrerem modificações, houve a criação de modelos de proteção e atendimento especial para os casos de violência doméstica, procurando proteger as mulheres agredidas antes de danos maiores, muitas vezes irreversíveis, incluindo aí o feminicídio, também incluído no Código Penal, como circunstância qualificadora do crime de homicídio, e incluí -lo no rol dos crimes hediondos.
A lei Maria da Penha busca realizar, como indicado no seu pórtico, a Convenção sobre a Eliminação de Todas as Formas de Discriminação contra as Mulheres e a Convenção Interamericana para Prevenir, Punir e Erradicar a Violência contra a Mulher.
No seu texto, várias são as formas de violência doméstica e familiar contra a mulher, identificadas como sendo, entre outras : “ a violência física, entendida como qualquer conduta que ofenda sua integridade ou saúde corporal ; a violência psicológica, entendida como qualquer conduta que lhe cause dano emocional e diminuição da autoestima ou que lhe prejudique e perturbe o pleno desenvolvimento ou que vise a degradar ou controlar suas ações, comportamentos, crenças e decisões, mediante ameaça, constrangimento, humilhação, manipulação, isolamento, vigilância constante, perseguição contumaz, insulto, chantagem, ridicularização, exploração e limitação do direito de ir e vir ou qualquer outro meio que lhe cause prejuízo à saúde psicológica e à autodeterminação ; a violência sexual, entendida como qualquer conduta que a constranja a presenciar, a manter ou a participar de relação sexual não desejada, mediante intimidação, ameaça, coação ou uso da força ; que a induza a comercializar ou a utilizar, de qualquer modo, a sua sexualidade, que a impeça de usar qualquer método contraceptivo ou que a force ao matrimônio, à gravidez, ao aborto ou à prostituição, mediante coação, chantagem, suborno ou manipulação ; ou que limite ou anule o exercício de seus direitos sexuais e reprodutivos ; a violência patrimonial, entendida como qualquer conduta que configure retenção, subtração, destruição parcial ou total de seus objetos, instrumentos de trabalho, documentos pessoais, bens, valores e direitos ou recursos econômicos, incluindo os destinados a satisfazer suas necessidades ; a violência moral, entendida como qualquer conduta que configure calúnia, difamação ou injúria. ” ( Conforme texto extraído do artigo sétimo da referida lei ).
Também se buscam medidas de integração entre os vários órgãos públicos, abrangendo repressão, proteção, conscientização, bem como utilização de equipes multidisciplinares, tendo em vista os graves danos, de variados matizes, que redundam das práticas agressivas.
Nessa linha, também existe na lei a previsão de locais de atendimento específico ( centros de atendimento integral e multidisciplinar para mulheres e respectivos dependentes em situação de violência doméstica e familiar ; casas-abrigo para mulheres e respectivos dependentes menores em situação de violência doméstica e familiar ; delegacias, núcleos de defensoria pública, serviços de saúde e centros de perícia médico-legal especializados no atendimento à mulher em situação de violência doméstica e familiar ; programas e campanhas de enfrentamento da violência doméstica e familiar ; centros de educação e de reabilitação para os agressores ).
O Poder Judiciário, também, deve especializar o seu atendimento para aplicar as medidas de restrição ao agressor de forma mais eficiente.
A linguagem utilizada na questão foi clara, indicando a grave violação dos direitos das mulheres injustamente agredidas e a luta pela busca de sanções adequadas aos agressores e de proteção pessoal e para os de sua família.
Ademais, o tema não está circunscrito ao Brasil, mas o país possui um número muito expressivo de atos que compõem esse cenário.
Por outro lado, o texto demanda uma reflexão sobre a realidade, e abarca diversos segmentos da sociedade, pois as situações de agressão não se restringem às camadas mais populares.
Como se trata de texto contendo uma ampla gama de possibilidades, é natural que ocorram interpretações diversas das esperadas.
O padrão de resposta identifica a projeção do elaborador na resposta a ser apresentada pelo estudante.
Está em linguagem correta e indica possíveis interpretações do texto.
Suas opções foram bastante amplas, abrangendo a maior parte das respostas encontradas.
No que concerne ao grau de profundidade, o padrão de resposta indicou a necessidade de conhecimento de lei especial que regula a violência doméstica, o que exige ou qualificação da área de Direito ou que os currículos dos cursos dos participantes desse ciclo do Enade tenham oferecido alguma disciplina introdutória da área jurídica.
Ao analisar as provas realizadas, constatamos que, na primeira parte, houve maior facilidade em responder aos itens indicados no padrão de resposta, pois atinentes, não só à formação profissional, mas, também, ao conhecimento da realidade, pois, infelizmente, a violência doméstica está disseminada na sociedade e nos vários pontos do país.
Mas, verificou -se que as respostas vieram, no mais das vezes, desorganizadas, dificultando a correção.
Apesar disso, restou claro que a maioria esmagadora das respostas confluiu para as indicações do padrão apresentado, sendo relevante indicar que o tema remeteu para a luta em relação à igualdade de gênero, temática candente em especial para as concluintes do sexo feminino.
Diversos textos se apresentaram de forma concatenada, fundamentada e, muitas vezes, lamentando que, nos dias atuais, ainda se discuta a necessidade de proteger as mulheres dessa espécie de violência.
Em muitos casos, também, os concluintes fizeram um liame entre o acesso à educação e a melhora das relações entre homens e mulheres, ao permitir que a conclusão de maior grau de educação formal conduzisse à diminuição do índice de violência contra a mulher.
Foram identificadas respostas consoantes com o padrão apresentado, indicando os itens que nele constam em relação aos reflexos pessoais, na família e na sociedade.
As respostas dos estudantes raramente faziam referência à organização protetiva instituída pela Lei Maria da Penha, como esperado pelo padrão de resposta.
A respeito de tal aspecto o padrão indicava como uma das possibilidades de respostas ( item ‘ b ’ ) : a criação de casas de refúgio ou casa-abrigo, para acolher mulheres vítimas de violência doméstica ; reforço às Delegacias de Atendimento à Mulher ; inclusão da vítima em programas assistenciais do governo, programas de proteção à vítima e à testemunha, transferência de local de trabalho, se a vítima for servidora pública.
Como paradigma, indicaram -se as Delegacias da Mulher que são organizadas, com a quase integralidade dos seus membros de policiais do sexo feminino, e especializadas no atendimento a mulheres submetidas a situações de violência.
Essa preocupação gerencial tem por escopo bloquear os preconceitos que acompanham essas dramáticas histórias de violência, o que gera uma incompreensão e uma ineficiência na rede de proteção.
As respostas demonstram que os concluintes têm ciência do problema, da sua gravidade, do vínculo com a questão de gênero e com a submissão da mulher como efeito da sociedade patriarcal, ainda que tratado de forma superficial, com base na mídia e, algumas vezes, com apoio de experiências pessoais.
Identificou -se uma generalizada compreensão de que o problema está sendo enfrentado e que ainda existe um longo caminho no percurso para a sua resolução. 3.3.1.5 - Análise de Língua Portuguesa das Questões Discursivas do Componente de Formação Geral Os dados de Farmácia, obtidos a partir das respostas às questões discursivas do Componente de Formação Geral, no que tange à Língua Portuguesa, encontram -se na Tabela 3.16 e no Gráfico 3.9.
Nesse aspecto, os alunos, de todo o Brasil, tiveram Média 59,6.
A maior Média com respeito à Língua Portuguesa foi obtida na região Sudeste ( 61,3 ), e a menor, na região Norte ( 54,0 ).
Quanto à variabilidade das notas, o Desvio padrão de todo o Brasil foi 20,7.
O menor Desvio padrão foi obtido na região Sudeste ( 19,5 ) e o maior Desvio padrão foi obtido na região Norte ( 22,6 ).
A Mediana das notas de Língua Portuguesa foi 65,0 para o Brasil como um todo, a mesma verificada nas regiões Nordeste e Centro-Oeste.
A maior Mediana com respeito à Língua Portuguesa foi obtida nas regiões Sudeste e Sul ( 67,5 ), e a menor, na região Norte ( 60,0 ).
A nota Máxima para todo o Brasil foi de 92,5, com, pelo menos, um aluno tirando essa nota nas regiões Sudeste e Centro-Oeste.
Além disso, a nota Mínima foi zero em todas as regiões do país.
O Gráfico 3.9 mostra a distribuição das notas de Língua Portuguesa do Componente de Formação Geral.
Observa -se que a maior frequência ( quase 35 \% ) corresponde aos alunos que obtiveram nota no intervalo ( 60 ; 70 ].
Destacam -se, também, os alunos que deixaram ambas as questões em branco, representando aproximadamente 7 \% do total. 3.3.1.6 - Comentários sobre a correção das respostas de Formação Geral com respeito à Língua Portuguesa Na avaliação do desempenho linguístico dos participantes foram consideradas quatro competências distintas, de modo a permitir um mapeamento detalhado do domínio dos recursos disponíveis na Língua Portuguesa para a comunicação escrita formal.
Nessa situação comunicativa, as exigências para o desenvolvimento desse tipo de texto dizem respeito : à adequação da seleção vocabular, ao desenvolvimento do conteúdo, à estruturação sintática dos períodos, à organização lógica das ideias, à utilização de procedimentos de encadeamento textual e de referenciação, à obediência às exigências morfossintáticas próprias da modalidade escrita da norma-padrão, ao respeito às regras ortográficas e às regras de acentuação gráfica.
A avaliação foi realizada considerando -se os textos das respostas às duas questões de Formação Geral e os resultados obtidos em cada uma delas, como seria de se esperar, foram semelhantes.
As competências avaliadas podem ser sintetizadas nos aspectos descritos a seguir.
Aspectos textuais - estruturação textual condizente com o gênero discursivo e o modo de organização textual expositivo/argumentativo adequado ao gênero – essa competência envolve : a estruturação sintática condizente com o padrão da modalidade escrita formal da Língua Portuguesa de modo a garantir a clareza necessária ; a distribuição do conteúdo do texto em parágrafos, de modo a garantir a sua organização temática ; a utilização de operadores discursivos que contribuam para a progressão temática do texto, estabelecendo relações lógicas entre as ideias apresentadas, tanto do ponto de vista intrafrasal, como do interfrasal ; a utilização de procedimentos de referenciação lexical e pronominal que permitam a retomada de referentes textuais ; o respeito às regras de pontuação como fator de estruturação do período.
Aspectos morfossintáticos - domínio dos diferentes aspectos morfossintáticos próprios da modalidade escrita formal da norma-padrão da Língua Portuguesa – essa competência envolve : a concordância nominal, a concordância verbal, a regência nominal, a regência verbal, a flexão nominal, a flexão verbal, a correlação entre os tempos verbais, a colocação pronominal e a utilização de sinais de pontuação que contribuam para a organização lógica da frase e do texto.
Aspectos vocabulares - seleção vocabular adequada à modalidade escrita formal da Língua Portuguesa, exigida pela situação comunicativa – essa competência envolve a precisão na seleção/utilização do vocabulário relacionado à temática solicitada pela questão ; a ausência de marcas da oralidade, como termos de sentido muito genérico ( “ coisa ”, “ negócio ”, “ você ” ) e termos de registros mais informais ( como gírias, jargões, frases feitas, ditados populares, termos regionais ).
Assim, espera -se que o participante respeite a adequação vocabular não usando gírias ou expressões coloquiais, evite repetição desnecessária de palavras e utilize um vocabulário mais formal, como solicitado por um texto dissertativo.
Aspectos ortográficos - respeito às convenções ortográficas da norma-padrão da Língua Portuguesa – essa competência envolve o domínio das regras de acentuação gráfica e da grafia padrão das palavras ( com ausência de abreviaturas próprias da linguagem da internet ), de acordo com as convenções estabelecidas pela legislação em vigor e consubstanciadas no Vocabulário Ortográfico da Língua Portuguesa, editado pela Academia Brasileira de Letras. ( Vale notar que, nessa edição de 2016, não será aceita a legislação anterior, no caso das regras relativas ao uso do hífen e da acentuação gráfica ).
A escolha dessas competências para subsidiar o processo de avaliação apoia -se na concepção de que, no desempenho dos graduandos, a modalidade escrita tem apresentado uma intensa simplificação, originada no padrão da modalidade oral da Língua Portuguesa.
No caso do texto de base dissertativa, inscrito em um registro formal, a distância entre as duas modalidades é ainda maior, o que provoca situações de hipercorreção ( desvios provocados pela incorporação indevida de uma regra da norma-padrão ) e de truncamentos sintáticos ( estruturas frasais incompreensíveis devido à complexidade sintática própria da modalidade escrita ).
Para efeito de pontuação, o critério aprovado pelo Inep, e já utilizado na avaliação do desempenho linguístico nas edições do Enade desde 2013, reúne as competências acima descritas em três grupos da seguinte forma : aspectos textuais ( 40 \% ), aspectos morfossintáticos e vocabulares ( 40 \% ) e aspectos ortográficos ( 20 \% ).
A seguir, analisam -se qualitativamente os resultados da avaliação de cada um dos aspectos avaliados nos textos de resposta das questões discursivas 1 e 2 de Formação Geral.
Aspectos textuais Esta competência é a que se revelou como a mais problemática entre os participantes, porque são muitos os desvios observados.
Em relação à estrutura formal do texto produzido pelos participantes, observaram -se relações linguísticas quase agramaticais, como as estabelecidas pela sequência de gerúndios sem o apoio de um ponto de partida para a organização das informações gramaticais e semânticas, ou seja, sem oração principal.
Observou -se que uma grande parte dos participantes não organizou as ideias em parágrafos, talvez devido ao pequeno número de linhas disponibilizadas para a resposta da questão ou, quem sabe, pela suposição de que não seria necessária essa divisão por não se tratar de um texto no modelo de uma redação dissertativo-argumentativa, como solicitado nos vestibulares.
Quanto à utilização dos mecanismos de referenciação, deve -se destacar a ocorrência, em uma parte dos textos, de repetições de palavras ou expressões sem a utilização de termos sinônimos ou pronomes, como seria adequado.
Outro aspecto relevante a ser destacado é a quase total ausência de operadores argumentativos, tanto intrafrasais como interfrasais, repercutindo uma tendência atual da mídia escrita.
A respeito do uso dos sinais de pontuação, observaram -se vários desvios nos textos analisados, com a ocorrência, inclusive, de parágrafos sem marca interna de pontuação para separar os períodos.
São os seguintes os tipos de problemas encontrados, que se caracterizam como desvios ao estabelecido pela norma-padrão : a ) uso inadequado de vírgula : • ocorrência entre o sujeito e o predicado, o que é considerado inadequado por interromper a sequência natural da frase ; • ocorrência de apenas uma das vírgulas para intercalar uma palavra, uma expressão ou uma oração encaixada, quando o correto seria a demarcação com uma vírgula antes e uma depois do termo encaixado na frase ; • ocorrência da vírgula no lugar do ponto para separar ideias que constituem períodos distintos ; b ) uso inadequado de ponto e vírgula para separar elementos que deveriam ser separados por vírgula, por estarem dentro do mesmo período ; c ) ausência de ponto final para separar períodos.
Vale observar que não foi penalizada a ausência de vírgula para destacar locuções ou adjuntos adverbiais de pequena extensão deslocados de posição na frase, por ser um uso opcional.
Os problemas relativos à estruturação textual geram, portanto, os seguintes tipos de desvios : • sequência justaposta de ideias sem encaixamentos sintáticos ; • redução de estruturas subordinadas, ao lado do aumento na frequência de estruturas coordenadas e absolutas ; • redução no uso de conectores para expressar relações lógicas essenciais à construção do texto ; • emprego inadequado de operadores que não estabelecem relações lógicas coerentes entre ideias do texto ; • repetição de termos sem a utilização de procedimentos mais sofisticados de substituição ( hiperonímias, hiponímias, nominalizações, expressões metafóricas ) ; • frases fragmentadas que comprometem a estrutura lógico-gramatical ; • frases formadas apenas por oração subordinada, sem oração principal.
Aspectos morfossintáticos e vocabulares Os resultados relativos aos aspectos morfossintáticos ( concordância, regência, colocação de pronomes, flexão nominal e verbal ) evidenciam desvios recorrentes no desempenho dos participantes, relacionados ao seu processo de formação e, principalmente, a hábitos da oralidade, marcados pela informalidade.
Apesar de exemplificarem processos de mudança em curso na Língua Portuguesa, esses desvios foram penalizados por não atenderem às exigências da norma-padrão relativas ao texto escrito em situação formal, como deveria ser o texto dissertativo solicitado no Enade.
O desvio mais frequente, em relação à regência, é a falta do sinal indicativo da crase – isso revela que o usuário não tem consciência de que, sob a forma do termo “ a ”, existe a presença de uma preposição “ a ”, exigida pela regência do termo anterior, existe a presença de uma preposição “ a ”, exigida pela regência do termo anterior, combinada a um artigo definido ( por exemplo : “ Devido as grandes guerras que estão acontecendo... ” ) Outros problemas relacionados à regência verbal e à nominal foram identificados em alguns textos dos participantes : ( a ) a ausência de preposição ou seu uso indevido antes de pronome relativo ( por exemplo : “ a região que eles moram não oferece condições de segurança ” ) ; ( b ) o uso indevido do pronome relativo ( por exemplo : “ não pensam na cidade para que eles vão ” ) ; ( c ) utilização de pronome “ onde ” deslocado de sua função locativa ).
Esses processos refletem hábitos da modalidade oral da língua, em situações de registro informal.
Apesar da possibilidade de que essa alteração de regência se generalize no padrão escrito da Língua Portuguesa, como já está ocorrendo até em textos jornalísticos, a ausência da preposição foi penalizada neste processo de avaliação.
Outro desvio observado nos textos diz respeito aos processos de concordância verbal e de concordância nominal.
Quanto à concordância de número, além das ocorrências já conhecidas de falta da marca de plural em verbos ou adjetivos, em função da ocorrência de substantivo plural na função de sujeito, observaram -se dois tipos de desvio já registrados nas pesquisas linguísticas : ( a ) a ausência de marca de plural com sujeito posposto ( por exemplo : “ Foi expulso do país os imigrantes que entraram sem autorização da polícia ” ) ; ( b ) o uso indevido da marca de plural para estabelecer concordância com o plural da locução adjetiva, apesar de o substantivo que funciona como núcleo do sintagma nominal estar no singular ( por exemplo : “ A proteção das mulheres que sofrem violência foram garantidas pela Lei Maria da Penha ” ).
Uma ocorrência generalizada foi a ausência de acento circunflexo na forma plural do presente do indicativo do verbo “ ter ”, que foi considerada como um desvio na concordância verbal e não meramente de acentuação gráfica.
Quanto à concordância de gênero, alguns casos foram observados, normalmente no âmbito de sintagmas nominais longos, em que o adjetivo está afastado do substantivo.
Quanto à questão da colocação pronominal, foram poucos os desvios observados.
Apesar de serem aspectos relacionados à oralidade, observou -se que a maioria dos participantes já incorporou regras como a não introdução da frase por um pronome oblíquo e a próclise na presença de um termo atrator.
Não se adotou, entretanto, o padrão excessivamente formal descrito pelas gramáticas normativas em relação à posição do pronome oblíquo em locuções verbais, já que esse uso está muito distante da prática cotidiana, até em textos mais formais.
Quanto aos aspectos vocabulares ( precisão semântica do vocabulário utilizado e adequação ao registro formal exigido pelo tipo de texto solicitado no enunciado ), podem ser destacados os seguintes tipos de desvios : expressões da oralidade, seleção vocabular incompatível com o contexto, gerando situações de falta de inteligibilidade ; falta de domínio de vocabulário mais abstrato e de maior complexidade, essencial ao desenvolvimento do texto de base dissertativa.
Não se observou presença de palavras de baixa calão, ao contrário do exame de 2015, em que o tema “ funk ” gerou a utilização desse tipo de vocabulário.
Várias marcas de oralidade foram identificadas, embora não em alta frequência : o uso do pronome relativo “ onde ” como relativo universal, falta de artigo definido antes de substantivo, repetição de palavras por falta de vocabulário, reduções como “ tá ”, “ pra ”, “ pro ”, “ prum ”, expressões informais.
Aspectos ortográficos : O desempenho dos participantes revelou uma diferença muito grande nos dois aspectos de caráter ortográfico analisados neste item : grafia e acentuação.
Observou -se a ocorrência de baixo índice de desvios na grafia das palavras e grande índice de desvios no uso dos sinais de acentuação gráfica.
Em vários casos, ocorre ausência completa de acentuação.
Os resultados revelam que existe uma tendência crescente entre os universitários brasileiros de eliminação da acentuação gráfica, talvez motivada pelos hábitos relacionados às redes sociais ou devido a um conhecimento ainda superficial sobre o Acordo Ortográfico da Língua Portuguesa de 1990.
Os casos mais sistemáticos de eliminação do acento indicador da sílaba tônica são : • palavras proparoxítonas ( por exemplo : “ proximos ”, “ politicos ”, “ publica ”, “ fisica ”, “ psicologica ” ) ; • palavras paroxítonas terminadas em ditongo crescente ( por exemplo : “ necessario ”, “ noticiarios ”, “ individuo ”, “ dependencia ”, “ varios ”, “ propria ”, “ violencia ”, “ ocorrencia ”, “ suicidio ” ) ; • palavras oxítonas ( por exemplo : “ ate ”, “ tambem ”, “ mes ” ) ; • palavras paroxítonas terminadas em ditongo nasal ( por exemplo : “ orgão ”, “ orfão ” ).
Quanto ao domínio das convenções relativas à grafia das palavras, pode -se concluir que os desvios estão majoritariamente relacionados a questões relativas à oralidade : • hipercorreção pela escolha de “ e ” no lugar de “ i ” ou de “ i ” no lugar de “ e ”, por influência de hábitos da oralidade : por exemplo, “ inviado ” ( no lugar de “ enviado ” ), “ entervenção ” ( no lugar de “ intervenção ” ), “ empedimento ” ( no lugar de “ impedimento ) ; • eliminação do “ r ” marcador do infinitivo verbal : por exemplo, “ esta ” ( no lugar de “ estar ” ), “ estuda ” ( no lugar de “ estudar ” ), “ migra ” ( no lugar de “ migrar ” ) ; • utilização de “ s ” no lugar de “ c ” : por exemplo, “ mensionado ” ( no lugar de “ mencionado ” ) ; “ insentivou ” ( no lugar de “ incentivou ” ) ; • utilização de “ s ” no lugar de “ x ” : por exemplo, “ espectativa ” ( no lugar de “ expectativa ”, devido à confusão com o verbo “ esperar ” ) ; e a situação contrária, “ x ” no lugar de “ s ”, “ expantoso ” ( no lugar de “ espantoso ” ) ; • utilização de “ ss ” no lugar de “ ç ” : por exemplo, “ gerassões ” ( no lugar de “ gerações ” ) ; • uso de “ am ” no lugar de “ ão ” na forma de terceira pessoa do singular do presente e do pretérito perfeito do indicativo : por exemplo, “ estam ” ( no lugar de “ estão ” ) ; • utilização de “ l ” no lugar de “ u ” em final de sílaba : por exemplo, “ tralma ” ( no lugar de “ trauma ” ) ; • outros casos de desvios de grafia relacionados à variação diastrática podem ser observados em “ estrupo ” ( no lugar de “ estupro ” ), “ subjulgadas ” ( no lugar de “ subjugadas ” ), “ precoseito ” ( no lugar de “ preconceito ” ), “ soubre ” ( no lugar de “ soube ” ), “ vecendo ” ( no lugar de “ vencendo ” ), “ indesencia ” ( no lugar de “ indecência ” ), “ apolojia ” ( no lugar de “ apologia ” ), “ fulga ” ( no lugar de “ fuga ” ). • inadequação no uso da maiúscula para destacar determinadas palavras-chave do texto, como “ Violência ”, “ Brasileiros ”.
Destaque -se, também, a ocorrência de participantes que grafam os textos inteiramente em caixa alta.
Vale observar que, ao contrário do esperado, não apareceram abreviaturas próprias do “ internetês ”, relacionadas ao uso de redes sociais e de e-mails.
Sua influência, portanto, parece se restringir à ausência dos acentos gráficos.
Em relação aos aspectos textuais, observa -se dificuldade na estruturação textual, o que resulta em truncamentos sintáticos e prejuízo da compreensão do texto, e desrespeito a regras essenciais de pontuação, principalmente pela dificuldade em identificar as relações lógicas que se estabelecem entre as unidades frasais.
Em relação aos aspectos morfossintáticos, domínio das regras de concordância nominal e verbal, mas desrespeito às regras de regência nominal e verbal, principalmente no que diz respeito ao uso do sinal indicativo da crase.
Em relação aos aspectos vocabulares, reduzida utilização de palavras de uso mais formal e abstrato, adequado à construção do texto de base dissertativa.
Em relação às convenções ortográficas, domínio da grafia das palavras, mas falta de atendimento à maioria das regras de acentuação gráfica. 3.3.2 - Componente de Conhecimento Específico Na parte da prova relativa às questões discursivas do Componente de Conhecimento Específico ( Tabela 3.17 ), observa -se que a Média foi um pouco mais alta do que para as questões discursivas do Componente de Formação Geral.
Enquanto no Componente de Formação Geral a Média para estudantes de Farmácia de todo o Brasil foi 54,8, na parte de Conhecimento Específico a Média foi 55,2.
A maior Média deste componente foi obtida pelos estudantes da região Sul ( 57,8 ), e a menor, pelos da região Norte ( 45,3 ).
Quanto à variabilidade das notas, o Desvio padrão de todo o Brasil foi 19,3.
O maior Desvio padrão foi encontrado na região Norte ( 20,7 ), e o menor, na região Sul ( 18,0 ).
A maior nota Máxima, 100, foi obtida em duas regiões : Nordeste e Sudeste.
Nas demais regiões a nota Máxima foi : 90,0 na região Norte, 93,3 na região Sul e 95,0 na região Centro-Oeste.
Além disso, a nota Mínima ( 0,0 ) foi obtida por pelo menos um aluno em todas as regiões do Brasil.
A Mediana do Brasil como um todo foi 56,7.
A maior Mediana foi obtida nas regiões Sudeste e Sul ( 60,0 ), e a menor, na região Norte ( 46,7 ).
O Gráfico 3.10 representa a distribuição das notas nas questões discursivas no Componente de Conhecimento Específico.
A moda desta distribuição ocorre no intervalo ( 60 ; 70 ], mesmo intervalo modal da distribuição das notas das questões discursivas de Formação Geral.
Observa -se, ainda, que essa distribuição é mais espalhada do que a das notas das discursivas de Formação Geral ( Gráfico 3.6 ).
A análise de cada uma destas questões será feita a seguir. 3.3.2.1 - Análise da Questão Discursiva 3 do Componente de Conhecimento Específico Na questão 3, cujos resultados aferidos encontram -se descritos na Tabela 3.18, a Média dos estudantes de todo o Brasil foi 64,4.
A menor Média nessa questão foi obtida pelos alunos da região Norte ( 56,2 ), enquanto a maior Média foi obtida na região Sul ( 67,4 ).
Quanto à variabilidade das notas, o Desvio padrão de todo o Brasil foi 24,3.
O maior Desvio padrão foi obtido na região Norte ( 26,6 ), enquanto o menor foi obtido na região Sul ( 22,8 ).
A nota Máxima, 100,0 pontos, foi alcançada por, pelo menos, um aluno de todas as regiões.
A Mediana do Brasil como um todo foi 70,0, a mesma para todas as regiões, exceto a da região Norte onde a Mediana foi 60,0.
A nota Mínima ( 0,0 ) foi obtida por, pelo menos, um estudante em todas as regiões do Brasil.
O Gráfico 3.11 mostra a distribuição das notas na questão discursiva 3, do Componente de Conhecimento Específico da área de Farmácia.
Essa distribuição apresenta moda no intervalo ( 70 ; 80 ], com maior concentração nas notas mais altas. 3.3.2.2 - Comentários sobre as respostas à Questão Discursiva 3 O enunciado da questão atende às exigências básicas das Diretrizes Curriculares da área de Farmácia.
A questão abordava um tema interessante e extremamente pertinente à formação do farmacêutico, situação que poderia ser encontrada no dia a dia do profissional que atua em farmácias e drogarias.
Os conhecimentos necessários costumam ser bastante discutidos em diferentes disciplinas do currículo dos Cursos de Farmácia, e, ainda, no estágio obrigatório realizado pelo estudante.
A questão descreveu, com detalhes, as infrações encontradas pelos fiscais do Conselho Regional de Farmácia ( CRF ) em um estabelecimento farmacêutico, assim como trouxe as considerações do profissional, quanto às suas infrações, quando submetido à audiência junto à comissão de ética do CRF.
O comando da questão se subdividia em três itens.
O nível de profundidade exigido pela questão pode ser considerado mediano, especialmente nos seus itens ‘ a ’ e ‘ b ’.
O item ‘ c ’ exigia do estudante um conhecimento mais específico da legislação – o tipo de sanção disciplinar imputado para as referidas infrações.
As respostas solicitadas em cada item, de maneira geral, eram pontuais, o que favoreceu aos estudantes serem, muitas vezes, sucintos, e terem tranquilidade quanto ao tempo necessário de resposta para cada item.
O espaço oferecido para a redação da resposta foi suficiente para a grande maioria dos estudantes.
O item ‘ a ’ solicitava ao estudante que citasse três infrações cometidas pelo farmacêutico e a grande maioria respondeu a esse item parafraseando o enunciado da questão, tal como foi apresentado.
Nesse item, as infrações mais apontadas foram : deixar o armário de medicamentos ( controlados ) aberto ; medicamentos com prazo de validade vencido ; agulhas de aplicação de injetáveis descartadas incorretamente e submetendo a população ao risco de contaminação.
Muitos estudantes, também, apontaram como infração a ausência do farmacêutico no estabelecimento.
Mas, raríssimos foram os casos em que descreveram que deve ser feita uma comunicação do afastamento temporário ao CRF ( por escrito ), quando não há outro farmacêutico que o substitua.
Outra resposta que foi muito apontada pelos estudantes foi mencionar que “ a chave deve ficar somente com o farmacêutico ”, mostrando saberam que não é permitido delegar a outros profissionais atos ou atribuições exclusivas da profissão farmacêutica, conforme estabelecido no padrão de resposta.
Um número considerável de respostas como negligência, irresponsabilidade e falta de ética também foram encontradas. 3.3.2.3 - Análise da Questão Discursiva 4 do Componente de Conhecimento Específico A Tabela 3.19 contém as informações relativas à questão 4 do conjunto de questões do Componente de Conhecimento Específico.
O desempenho dos estudantes nessa questão foi inferior ao das questões 3 e superior ao da questão 5.
A Média geral No item ‘ b ’, o comando solicitava que o estudante apontasse três implicações dessas infrações na sociedade, no reconhecimento e na valorização da profissão pela comunidade.
Semelhante ao item anterior, grande parte dos estudantes descreveu como uma das implicações das infrações “ a desvalorização da profissão pela sociedade ”, frase que consta do comando do item.
Nesse item, as implicações mais apontadas foram : o “ uso irracional de medicamentos ” e “ uso de medicamentos vencidos, com possíveis consequências indesejáveis para a saúde ”.
Um elevado percentual de estudantes apontou como implicações “ deixar de prestar assistência farmacêutica ” ou “ atenção farmacêutica".
A leitura das respostas para esse item deixou claro que os estudantes não conseguem diferenciar os conceitos de atenção farmacêutica e de assistência farmacêutica.
No item ‘ c ’, entre as respostas esperadas, ' multa ' foi a sanção mais citada.
Houve um número considerável de estudantes que apontou que o farmacêutico perderia seu CRF, sem descrever que antes dessa perda há outras sanções passíveis de serem aplicadas dentro da legislação.
O enunciado da questão sugere que foi a primeira vez que o profissional foi convocado pelo comitê de ética, no entanto, não deixa isto em evidência, o que pode ter levado os estudantes a considerar a perda do CRF como resposta.
Um grande número de estudantes apenas escreveu “ aplicação de sanção ”, demonstrando a falta de conhecimento específico da legislação.
Quase a totalidade dos estudantes tentou responder à questão, diferentemente das outras duas questões discursivas, portanto, o percentual de respostas em branco foi o menor entre as três questões discursivas.
Tal fato demonstra que o tema abordado, muito pertinente, tem sido amplamente discutido com os estudantes em sua formação por todo o país.
As respostas apresentadas apontaram o crescimento do entendimento quanto à importância da área de assistência farmacêutica na formação profissional.
Por outro lado, as respostas também demonstraram que os estudantes, de maneira geral, devem ser estimulados a redigir respostas para questões discursivas durante a graduação, de modo a melhorar a capacidade de elaboração de textos. 3.3.2.3 - Análise da Questão Discursiva 4 do Componente de Conhecimento Específico A Tabela 3.19 contém as informações relativas à questão 4 do conjunto de questões do Componente de Conhecimento Específico.
O desempenho dos estudantes nessa questão foi inferior ao das questões 3 e superior ao da questão 5.
A Média geral do Brasil foi 56,0, sendo a menor Média registrada na região Norte ( 40,4 ), e a maior, na região Sul ( 64,6 ).
A nota Máxima ( 100,0 ) foi atingida por, pelo menos, um concluinte de todas as regiões.
A Mediana em todo o Brasil foi 60,0, o mesmo valor foi obtido nas regiões Nordeste, Sudeste e Centro-Oeste.
Nas demais regiões a Mediana foi : 40,0 na região Norte e 70,0 na região Sul.
As notas mínimas foram zero em todas as regiões.
O Gráfico 3.12 representa a distribuição de notas da questão discursiva 4, do Componente de Conhecimento Específico.
Essa distribuição tem moda no intervalo ( 70 ; 80 ] e apresenta máximo local no intervalo ( 30, 40 ].
Em torno de 8 \% dos alunos deixaram a resposta em branco. 3.3.2.4 - Comentários sobre as respostas à Questão Discursiva 4 A questão era pertinente em relação às Diretrizes Curriculares da área de Farmácia.
Como na questão 3, abordava tema relacionado à prática diária do profissional farmacêutico.
O enunciado descrevia o caso de um paciente que apresentava histórico familiar de trombose e também trazia os resultados de um exame laboratorial com parâmetros alterados.
Informou, ainda, que após consulta médica, o clínico prescreveu os seguintes medicamentos : sinvastatina 10 mg, hidroclorotiazida 50 mg e ácido acetilsalicílico ( AAS ) 100 mg, fatos pertinentes em relação ao quadro clínico descrito.
A questão apresentou total suficiência de informações e, sendo assim, os concluintes teriam plenas condições de respondê -la.
O comando da questão, no item ‘ a ’, solicitava que o estudante justificasse a terapia prescrita, relacionando cada fármaco ao histórico clínico e aos exames laboratoriais apresentados, o que está tecnicamente adequado.
No item ‘ b ’, o comando solicitava que, para cada fármaco prescrito, o estudante citasse dois exames laboratoriais de acompanhamento para o paciente.
Para este item, o padrão apontava Uma ampla possibilidade de opções de respostas que poderiam ser aceitas como corretas.
Cabe ressalvar que, para o acompanhamento do AAS, o padrão de respostas sugeria que os exames de tempo de coagulação sanguínea e contagem de plaquetas, procedimento corroborado pela grande maioria das solicitações médicas.
O nível de profundidade exigido pode ser considerado de mediano a difícil em relação às outras questões discursivas.
A questão não possibilitava interpretações diversas, nem o comando abria espaço para divagações.
Tal fato permitiu certa homogeneidade das respostas apresentadas.
Nas respostas, não foram poucos os alunos que descreveram o ácido acetilsalicílico ( AAS ) como anticoagulante.
Um fato marcante, no que tange à linguagem, foi que um número considerável de estudantes se refere ao efeito do AAS no sangue como"afinar o sangue".
Trata -se de linguagem completamente inadequada para um profissional da área de saúde, além de não ter nenhum significado no âmbito profissional ou científico.
Outro ponto importante a ser ressaltado foi o alto percentual de alunos que descreveram que a sinvastatina e o AAS podem ser avaliados por meio de um hemograma, demonstrando desconhecimento de quais parâmetros são avaliados nessa solicitação laboratorial.
Tal fato é extremamente preocupante, especialmente após a publicação da resolução nº 585, de 29 de agosto de 2013, que regulamenta as atribuições clínicas do farmacêutico e em cujos incisos XI a XIV ( artigo 7º, capítulo I ) está descrito que o farmacêutico poderá : solicitar exames laboratoriais, no âmbito de sua competência profissional, com a finalidade de monitorar os resultados da farmacoterapia ; avaliar resultados de exames clínico-laboratoriais do paciente, como instrumento para individualização da farmacoterapia ; monitorar níveis terapêuticos de medicamentos, por meio de dados de farmacocinética clínica ; determinar parâmetros bioquímicos e fisiológicos do paciente, para fins de acompanhamento da farmacoterapia e rastreamento em saúde.
Esta questão, entre as três discursivas, foi a que apresentou o maior percentual tanto de notas zero, quanto de respostas em branco.
O percentual de respostas consideradas fracas foi menor que 20 \%.
As respostas medianas foram as que apresentaram o maior percentual, pouco mais de 40 \%.
Em torno de 1/3 das respostas foram consideradas boas.
Pouco mais de 5 \% dos estudantes obtiveram nota 10,0 na questão.
As soluções apresentadas pelos estudantes mostraram algumas deficiências em conceitos importantes e fundamentais para a formação de um profissional farmacêutico, como, por exemplo, não estabelecer a diferença entre anticoagulante e antiagregante plaquetário.
Outro exemplo, que também acentua a preocupação com a formação na Área, é o desconhecimento dos estudantes sobre os parâmetros que podem ser obtidos em exames laboratoriais, como o hemograma.
Ficou evidente que muitos estudantes não estabeleceram a diferença entre eletrólitos e a presença de sódio e potássio. 3.3.2.5 - Análise da Questão Discursiva 5 do Componente de Conhecimento Específico A Tabela 3.20 contém as informações relativas à questão 5 do conjunto do Componente de Conhecimento Específico.
O desempenho dos estudantes de todo o Brasil nesta questão foi o mais baixo dentre as três questões discursivas desse componente.
A nota Média dos estudantes de todo o Brasil foi 45,1.
A maior Média foi registrada na região Sudeste ( 47,8 ), enquanto a menor Média foi registrada na região Norte ( 39,2 ).
Quanto à variabilidade das notas, o Desvio padrão dos alunos do Brasil, como um todo, foi 25,5.
Enquanto o maior desvio foi encontrado na região Norte ( 26,0 ), o menor foi encontrado na região Sul ( 24,7 ).
A Mediana para o Brasil foi 50,0, a mesma das regiões Nordeste, Sudeste e Centro-Oeste.
As regiões Norte e Sul apresentaram uma Mediana mais baixa ( 40,0 e 45,0, respectivamente ) que a do Brasil como um todo.
Para o conjunto de alunos de Farmácia do Brasil, a nota Máxima foi 100,0, e a nota Mínima ( 0,0 ).
Os mesmos limites foram encontrados em todas as regiões.
O Gráfico 3.13 representa a distribuição das notas da questão discursiva 5 do Componente de Conhecimento Específico.
Essa distribuição tem moda no intervalo ( 50 ; 60 ] e mais dois máximos locais em ( 10 ; 20 ] e nas respostas “ em branco ”. 3.3.2.6 - Comentários sobre as respostas à Questão Discursiva 5 A questão discursiva 5 abordou um tema importante, atual e muito frequente no cotidiano do atendimento farmacêutico em farmácias e drogarias, assim como em todas as atuações relacionadas à saúde – as epidemias de dengue, zika e chikungunya.
O enunciado era claro e de fácil compreensão, apontando a importância das ações de controle para estas arboviroses, uma vez que não há tratamento específico.
Sendo assim, a questão levanta a importância do conhecimento do farmacêutico na orientação aos usuários do sistema de atenção à saúde e solicitava que o estudante respondesse a três itens : a ) indicar um repelente químico utilizado por gestantes e crianças menores de 2 anos ( segundo a OMS ) ; b ) citar riscos da automedicação na vigência da dengue ; e c ) citar medidas preventivas para reduzir a incidência dessas arboviroses.
As respostas esperadas para os itens ‘ a ’ e ‘ b ’ eram muito pontuais, já que esses itens não possibilitavam muitas variações.
Já para o item ‘ c ’, o padrão de repostas era amplo o suficiente para abranger variações que foram apresentadas pelos estudantes.
Quanto ao nível de profundidade, a questão pode ser considerada de fácil ( itens ‘ b ’ e ‘ c ’ ) a mediana ( item ‘ a ’ ).
O primeiro item ( ‘ a ’ ) teve um baixíssimo índice de respostas e dentre elas poucas estavam corretas.
Muitos estudantes apontaram o uso do óleo de citronela e de produtos naturais como repelentes seguros para gestantes e crianças menores de 2 anos.
No entanto, o comando deixava claro que o repelente deveria seguir as recomendações da Organização Mundial da Saúde, indicando como mais seguros o IR3535 e a icaridina ( crianças de 6 meses a 2 anos ), que não apresentam restrições para gestantes.
No item ‘ b ’, em geral, os estudantes foram pontuais descrevendo, somente, mas corretamente, a ' hemorragia ' como um dos possíveis riscos da automedicação.
Alguns estudantes a correlacionaram ao uso do ácido acetilsalicílico.
No item ‘ c ’, as principais medidas preventivas citadas pelos estudantes foram o controle de locais com água parada ( vasos, pneus, etc. ), uso de telas, uso de roupas adequadas e uso de repelentes.
A questão apresentou um percentual aproximado de 25 \% de respostas consideradas fracas ( notas até 30 pontos inclusive ).
Destas, cerca de 2 \% foram de notas zero.
As respostas consideradas medianas ( notas entre 35 e 70 ) representaram mais de 50 \% do total, o maior percentual.
Menos de 10 \% dos estudantes obtiveram notas boas ( notas iguais ou superiores a 75 ), entre estas, aproximadamente 1 \% recebeu a nota máxima, sendo este o menor percentual de notas 100 entre as três questões.
O percentual de provas em branco ficou acima de 5 \%.
Quase a totalidade dos estudantes apresentou a resposta subdividida em itens, como solicitado pelo comando.
Como as respostas deveriam ser muito objetivas, não houve muita diversidade nas soluções apresentadas.
O desconhecimento de quais os repelentes químicos são recomendados para crianças de 6 meses a 2 anos e para gestantes foi notado com base nas respostas dos estudantes.
Pelo fato de que este tema tem sido, sistematicamente, assunto de várias campanhas lançadas pelo Ministério da Saúde para o controle do Aedes aegypti, torna -se relevante frisar que muitos estudantes, em suas respostas, não demonstraram conhecimento dessa informação de domínio público.

% matched lemmas: 
\end{textsample}
